% Clean pre-results source for the ICIIT 2027 manuscript.
% The default build never prints editorial notes or result placeholders:
%   TEXINPUTS=..: latexmk -pdf -interaction=nonstopmode -halt-on-error main.tex
% Build review.tex instead when collaborators need the internal result outline.
%
% TITLE DECISION LOG
% Recommended and currently selected:
%   Lexicographic MaxSAT for Home-Care Resource Allocation:
%   Encoding and Constraint Strengthening
% Rationale:
%   - "Lexicographic MaxSAT" foregrounds the objective policy and formalism.
%   - "Home-Care Resource Allocation" retains the application and relationship
%     to the original HCORAP study.
%   - The subtitle groups Totalizer, implied constraints, and symmetry breaking
%     without turning the title into a list or claiming a uniform speedup.
%
% Shorter alternatives, in decreasing order of preference:
%   1. Lexicographic MaxSAT for Home-Care Resource Allocation: Evaluating
%      Totalizer Encoding and Constraint Enhancements
%   2. From Weighted to Lexicographic MaxSAT for Home-Care Resource Allocation
%   3. A Reproducible Evaluation of Lexicographic MaxSAT for Home-Care Resource
%      Allocation
% Avoid before data freeze: "faster", "superior", "robust", "Pareto-optimal",
% "routing-aware", and "state of the art".  The implemented model and frozen
% campaign do not currently support those title-level claims.

% The locally supplied acmart v1.90a loads hyperxmp before hyperref.  Current
% hyperxmp releases reject that legacy order, so preload hyperref without
% modifying the conference class.  Remove this compatibility line if a newer
% official template is supplied.
\RequirePackage[bookmarksnumbered,unicode]{hyperref}
\documentclass[sigconf,nonacm]{acmart}

\usepackage{booktabs}
\usepackage{etoolbox}
\usepackage{placeins}
\usepackage{tabularx}

% The default manuscript is visually clean.  review.tex defines
% \HCORAPInternalReview before loading this file and therefore exposes the
% editorial outline in a separate PDF.
\newif\ifinternaloutline
\ifdefined\HCORAPInternalReview
  \internaloutlinetrue
\else
  \internaloutlinefalse
\fi

% File existence alone is not evidence of a validated data freeze.  The marker
% is generated only by experiments/freeze_manuscript_bundle.py after it checks
% the primary analysis, every enabled corrected-v2 branch, cross-paradigm
% validation, screening decision, clean publication commit, and all three
% manuscript fragments.
\newif\iffrozenresults
\frozenresultsfalse
\IfFileExists{generated/freeze-manifest.tex}{%
  \input{generated/freeze-manifest}%
  \ifdefstring{\HCORAPFrozenValidationStatus}{VALID}{%
    \ifcsdef{HCORAPFrozenEvidenceDigest}{%
      \ifcsdef{HCORAPFrozenSourceCommit}{%
        \IfFileExists{generated/abstract-findings.tex}{%
          \IfFileExists{generated/results.tex}{%
            \IfFileExists{generated/conclusion.tex}{%
              \frozenresultstrue
            }{}%
          }{}%
        }{}%
      }{}%
    }{}%
  }{}%
}{}

% submission.tex activates a hard guard: a release build must never succeed
% without the complete, frozen evidence bundle.
\ifdefined\HCORAPSubmissionBuild
  \iffrozenresults\else
    \PackageError{hcorap-manuscript}{Submission build requires frozen results}{%
      Generate the three manuscript fragments from validated primary evidence,
      then run experiments/freeze_manuscript_bundle.py to issue the marker.}%
  \fi
\fi

\newcommand{\outlineblock}[1]{%
  \ifinternaloutline
    \par
    {\small\color{blue}\noindent\textbf{Outline note.} #1\par}
  \fi
}
% The argument identifies the expected frozen result (for example, n, x, or a
% generated visual) so that collaborators can distinguish pending artifacts.
\newcommand{\resultplaceholder}[1]{%
  \ifinternaloutline
    \textcolor{red}{[#1]}%
  \fi
}
\newcommand{\citationplaceholder}[1]{%
  \ifinternaloutline
    \textcolor{red}{[citation pending]}%
  \fi
}

% The conference has not supplied rights metadata yet.  Do not add DOI, ISBN,
% price, copyright, or proceedings metadata until the official instructions do.
\settopmatter{printacmref=false}
\renewcommand\footnotetextcopyrightpermission[1]{}

\begin{document}

\title[Lexicographic MaxSAT for HCORAP]{%
  Lexicographic MaxSAT for Home-Care Resource Allocation:
  Encoding and Constraint Strengthening}

% TODO by 2026-08-31: obtain all co-authors' confirmation of author order,
% corresponding author, ASCII spellings, and ORCIDs.  Keep these forms
% identical in the PDF, submission system, ORCID records, and artifact metadata.
\author{Tuyen Van Kieu}
\email{tuyenkv@vnu.edu.vn}

\author{Khanh Ngoc Do}
\email{2302061@vnu.edu.vn}

\author{Khanh Van To}
\email{khanhtv@vnu.edu.vn}
% Consecutive \author entries followed by one \affiliation form a shared
% affiliation group in acmart while retaining separate author metadata.
\affiliation{%
  \department{Faculty of Information Technology}
  \institution{VNU University of Engineering and Technology}
  \city{Hanoi}
  \country{Vietnam}}
\renewcommand{\shortauthors}{Kieu et al.}

\begin{abstract}
Home-care resource allocation assigns qualified caregivers to requested
services and time slots under coverage, availability, and workload constraints.
The existing MaxSAT formulation uses a weighted aggregate that can conceal
trade-offs among caregiver fragmentation, excess workload, and assignment
similarity.  We introduce LEX-COS, an exact staged policy that lexicographically
minimizes caregiver fragmentation and excess workload before maximizing
assignment similarity.  We combine it with bidirectional Totalizer encodings,
logically implied constraints, and symmetry breaking for detected slot and
service equivalences.  A pre-specified paired campaign separates
objective-policy effects from encoding interactions on the original and
corrected-v2 benchmarks.  It uses cumulative stage budgets, independent
solution verification, and optimal-objective checks against Gurobi and CPLEX.
\iffrozenresults
  % generated/abstract-findings.tex
% DATA PROVENANCE — MANDATORY NOTE:
%   Solver: EvalMaxSAT (EXPLORATORY). Primary campaign requires a fresh frozen run.
%   Replace with fresh publication-campaign results before submission.
%   Weighted: 250 instances, class 30_10_4. Commercial: 94 instances, class 30_15_4.
%   LEX-COS: 0 OPTIMUM (EvalMaxSAT timeout) — do NOT report as LEX-COS result.
Across the weighted $2^3$ factorial on 250 total instances (176 feasible), the Totalizer
encoding reports 163--163~solved (of~176) versus 158--161 for sorting networks,
with a median paired speedup of 1.12.  Gurobi and CPLEX agree on all 94~mutually
optimal weighted solutions; EvalMaxSAT agrees on all~19~overlapping instances.

\else
  \outlineblock{Insert two or three quantitative findings generated from the
  frozen tables: verified OPTIMUM counts, a paired objective change, and
  cross-framework objective agreement.  Do not use legacy or pilot values.}
\fi
\end{abstract}

\keywords{home health care scheduling, MaxSAT, lexicographic optimization,
Totalizer encoding, implied constraints, symmetry breaking}

\maketitle

\section{Introduction}
\label{sec:introduction}

Home-care planning must match requested services with qualified caregivers and
feasible time slots under coverage, availability, workload, and slot-conflict
constraints.  Schedule quality also depends on caregiver--service suitability,
continuity of care across related services, and excess workload.
We study the assignment-and-scheduling scope of the Home Care Optimal Resource
Allocation Problem (HCORAP): locations may inform assignment-suitability
rewards, but travel times, routes, and visit-to-visit transitions are not
decision variables.

Unceta et al.~\cite{UncetaEtAl2024} encode HCORAP as weighted partial MaxSAT.\@
A
weighted scalar score is useful when its coefficients faithfully encode the
intended exchange rates.  It does not, however, express a strict managerial
priority unless the weights are chosen to dominate every possible change in
lower-priority criteria.  Objective vectors with different CONT, OT, and SIM
components can consequently receive the same scalar score.  At the encoding
level, threshold variables used for workload and continuity can be represented by
different CNF cardinality encodings.  These representations preserve the
scheduling semantics but may differ materially in formula size, propagation,
and search behavior~\cite{BailleuxBoufkhad2003,AsinEtAl2011}.

This work makes four contributions.  First, it specifies an exact full-coverage
lexicographic policy, LEX-COS, and an alternative Overtime--Continuity--Similarity
order (LEX-OCS) for sensitivity analysis.  Second, it implements bidirectional
Totalizer reifications and two
families of logically implied constraints.  Third, it detects exact equivalence
classes of slots and services and adds symmetry-breaking constraints that retain
a canonical representative from every detected orbit, and hence at least one
optimal solution.  Fourth, it pre-specifies a paired evaluation on the original
and corrected-v2 benchmarks, with independent solution verification and
checks of optimal objective agreement against Gurobi and CPLEX.\@ We ask:
\emph{RQ1}, how do weighted and lexicographic policies differ in independently
recomputed objective vectors and performance in establishing optimality?
\emph{RQ2}, how does Totalizer encoding affect formula size, memory, and
performance in establishing optimality relative to sorting networks?
\emph{RQ3}, how do the encoding, implied-constraint, and symmetry-breaking
factors interact?

% RQ-to-evidence map (keep in source; express compactly in final prose):
% RQ1 weighted vs. LEX-COS/LEX-OCS -> paired objective vectors, solved, PAR-2.
% RQ2 sorting network vs. Totalizer -> size, RSS, solved, PAR-2, time to optimality.
% RQ3 implied x symmetry interactions -> paired 2x2x2 factorial analysis.
% Validation -> corrected-v2 evaluation set and three-backend objective agreement.

\section{Related Work}
\label{sec:related-work}

Home-care optimization spans assignment, scheduling, routing, and uncertainty.
Joint models can coordinate caregiver assignment with schedules and
routes~\cite{CappaneraScutella2015}, whereas stochastic assignment models address
uncertain service requirements~\cite{ErrarhoutEtAl2016}.  These richer scopes
motivate important extensions, but they are distinct from the deterministic,
time-slotted HCORAP studied here.  In Boolean optimization, weighted partial
MaxSAT has supported staff scheduling~\cite{DemirovicEtAl2019}, and Unceta et
al.~\cite{UncetaEtAl2024} give the MaxSAT formulation and benchmark on which our
study builds.

Lexicographic Boolean optimization and multi-objective MaxSAT are established
research areas~\cite{MarquesSilvaEtAl2011,JabsEtAl2024}.  Likewise, Totalizers
and cardinality networks provide alternative CNF representations of counting
constraints~\cite{BailleuxBoufkhad2003,BailleuxBoufkhad2004,AsinEtAl2011};
cardinality structure also matters inside MaxSAT
algorithms~\cite{MorgadoEtAl2014,JahrenAsin2018}.  Implied constraints can alter
propagation while preserving the feasible set, whereas symmetry breaking may
discard labeled solutions while retaining a representative of every detected
orbit~\cite{BofillEtAl2022,BogaertsEtAl2022}.  Our contribution is not to
introduce these general techniques, but to define their HCORAP semantics,
verify their safety conditions, and evaluate their interactions with a
domain-specific lexicographic policy.

\begin{table*}[t]
  \caption{Methodological scope of representative related-work streams.  A/S
  denotes assignment or scheduling; IC/SB denotes implied constraints or
  symmetry breaking; N/A means that cardinality encoding is not a study factor.}
  \label{tab:related-work}
  \centering
  \footnotesize
  \begin{tabularx}{\textwidth}{@{}Xccccccc@{}}
    \toprule
    Research stream & A/S & Routing & Uncertainty & Lex. & MaxSAT & Cardinality & IC/SB \\
    \midrule
    Joint home-care models~\cite{CappaneraScutella2015} & Yes & Yes & No & No & No & No & No \\
    Uncertain home-care assignment~\cite{ErrarhoutEtAl2016} & Yes & No & Yes & No & No & No & No \\
    MaxSAT staff scheduling~\cite{DemirovicEtAl2019} & Yes & No & No & No & Yes & Varies & No \\
    Original HCORAP~\cite{UncetaEtAl2024} & Yes & No & No & No & Yes & Sorting & No \\
    Boolean lexicographic and bi-objective MaxSAT~\cite{MarquesSilvaEtAl2011,JabsEtAl2024} & No & No & No & Yes & Yes & N/A & No \\
    Cardinality encodings in MaxSAT~\cite{BailleuxBoufkhad2003,BailleuxBoufkhad2004,MorgadoEtAl2014,JahrenAsin2018,AsinEtAl2011} & No & No & No & No & Yes & Yes & No \\
    Implied constraints and symmetry breaking~\cite{BofillEtAl2022,BogaertsEtAl2022} & No & No & No & No & Yes & No & Yes \\
    This work & Yes & No & No & Yes & Yes & Sorting/Totalizer & Yes \\
    \bottomrule
  \end{tabularx}
\end{table*}

\section{Problem Formulation and Objective Policies}
\label{sec:problem}

\subsection{HCORAP model and quality measures}

Let $\mathcal A$, $\mathcal S$, $\mathcal H$, $\mathcal U$, and $\mathcal Q$
denote caregivers, requested services, time slots, users, and continuity sets,
respectively.  A caregiver is called an \emph{agent} in the benchmark files,
and the benchmark field \texttt{SEQ} calls each continuity set a service
sequence even though no temporal ordering is modeled within it.  Each service
is assigned to one slot and contributes one workload unit.  A reward
$r_{as}\in\{0,\ldots,4\}$ measures the suitability of assigning caregiver $a$
to service $s$ and is the original HCORAP \emph{similarity score}; a zero reward
denotes a forbidden or unqualified assignment.
Availability, service-slot suitability, and qualification define the feasible
triples $E\subseteq\mathcal A\times\mathcal S\times\mathcal H$.  We create
$x_{ash}$ only for $(a,s,h)\in E$ and channel
$y_{as}\leftrightarrow\bigvee_{h:(a,s,h)\in E}x_{ash}$.

The hard model assigns every service exactly once.  At most one service may be
assigned to a caregiver in a slot, and at most one of a user's services may
occupy that slot.  If $L_a=\sum_s y_{as}$ is caregiver $a$'s workload, then
$L_a\le HN_a+HE_a$, where $HN_a$ and $HE_a$ are the regular and permitted
excess-workload limits in the same unit as assigned services.  The model does
not represent heterogeneous service durations, so it does not interpret these
units as clock hours.  Thus full coverage---every requested service is assigned
exactly once---is a hard requirement in every primary experiment; partial
schedules are not compared as policy outcomes.

Let $\mathcal S_u$ and $\mathcal S_q$ be the services of user $u$ and continuity
set $q$, respectively, and let
$n_q=|\{a:\exists s\in\mathcal S_q,\ y_{as}=1\}|$.  The measures are
\begin{align}
  \mathrm{SIM}  &= \sum_{a\in\mathcal A}\sum_{s\in\mathcal S} r_{as}y_{as}, \\
  \mathrm{CONT} &= \sum_{q\in\mathcal Q}\max(0,n_q-1), \\
  \mathrm{OT}   &= \sum_{a\in\mathcal A}\max(0,L_a-HN_a).
\end{align}
Here SIM is the assignment-suitability reward and is maximized; CONT counts the
additional distinct caregivers beyond the first in each continuity set and is
a caregiver-fragmentation penalty; and OT counts excess-workload units.  Both
CONT and OT are minimized.  Under full coverage, the stability reward in the
original HCORAP formulation is
\begin{equation}
  \mathrm{STAB}
  = \sum_{q\in\mathcal Q}(|\mathcal S_q|-n_q)
  = \sum_{q\in\mathcal Q}(|\mathcal S_q|-1)-\mathrm{CONT}.
  \label{eq:stability-continuity}
\end{equation}
Equation~\eqref{eq:stability-continuity} shows that maximizing the original
stability reward is equivalent to minimizing CONT.  For compatibility with the
released result schema, the
artifact fields remain named \texttt{similarity}, \texttt{continuity}, and
\texttt{overtime}; throughout this paper, however, CONT denotes the penalty,
not the desirable continuity-of-care level.

\subsection{Weighted baseline and lexicographic policies}

With $p=|P|$ denoting the non-negative magnitude of the per-instance signed
penalty field $P\leq0$, policy B0 maximizes
\begin{equation}
  \mathrm{SIM}-w_c\mathrm{CONT}-w_o\,p\,\mathrm{OT},
  \label{eq:weighted}
\end{equation}
where $w_c,w_o\in\mathbb Z_{\geq0}$ multiply CONT and OT; the primary experiments
use $(w_c,w_o)=(1,1)$.  The main policy solves
\begin{equation}
  \text{LEX-COS:}\qquad
  \min\mathrm{CONT}\ \longrightarrow\ \min\mathrm{OT}\
  \longrightarrow\ \max\mathrm{SIM}.
  \label{eq:lex-cos}
\end{equation}
LEX-OCS swaps the first two stages and is used only as an order-sensitivity
analysis.

Let $(f_1,f_2,f_3)=(\mathrm{CONT},\mathrm{OT},-\mathrm{SIM})$ and let $F_0$
be the schedules satisfying all hard constraints.  For $i=1,2,3$, the staged
solver constructs
\begin{equation}
  z_i=\min_{x\in F_{i-1}} f_i(x),\qquad
  F_i=\{x\in F_{i-1}:f_i(x)=z_i\}.
  \label{eq:lex-stages}
\end{equation}
Each stage rebuilds the WCNF over $F_{i-1}$, receives only the unused part of
one cumulative timeout, and passes its assignment to the independent verifier.
The next stage is entered only after EvalMaxSAT reports \textsc{Optimum}; the
next-stage WCNF preserves CONT and OT optima with upper pseudo-Boolean bounds
and the SIM optimum with a lower pseudo-Boolean bound.  Because each $z_i$ is
solver-established, the corresponding one-sided bound is equivalent to
$f_i(x)=z_i$ over $F_{i-1}$.  Thus $F_i$ contains precisely the schedules
attaining prefix $(z_1,\ldots,z_i)$, and a verified assignment in $F_3$ is
lexicographically optimal for the encoded model.  Any incomplete stage is a
policy timeout, not an optimum.  LEX-COS removes weight calibration for the
stated priority order; it does not dominate every stakeholder weighting.

\section{Encoding and Constraint Enhancements}
\label{sec:encoding}

\subsection{Sorting-network and Totalizer cardinality encodings}

The reified threshold vectors encode caregiver workload and the number of distinct
caregivers per continuity set.  The baseline uses a bidirectional sorting
network~\cite{AsinEtAl2011}.  The alternative uses the standard Totalizer tree
of Bailleux and Boufkhad~\cite{BailleuxBoufkhad2003}, completed with both merge
clause families from their full-CNF treatment~\cite{BailleuxBoufkhad2004}.
Thus, this is a full threshold reification rather than a truncated
forward-only at-most-$k$ encoding.  When child vectors $L$ and $R$, of lengths
$m$ and $n$, are merged into $O$, use $L_0=R_0=\top$ and
$L_{m+1}=R_{n+1}=\bot$ and add
\begin{align}
  \neg L_{\alpha}\lor\neg R_{\beta}\lor O_{\alpha+\beta},\qquad
  L_{\alpha+1}\lor R_{\beta+1}\lor\neg O_{\alpha+\beta+1},
\end{align}
for $0\leq\alpha\leq m$ and $0\leq\beta\leq n$; the first clause is included
when $\alpha+\beta\geq1$ and the second when $\alpha+\beta<m+n$.  Together they enforce
$O_k\leftrightarrow(\sum_j x_j\ge k)$.  Exact threshold semantics matter here
because these auxiliary literals occur directly in objectives and stage-fixing
equality bounds; a one-directional implication would not suffice.  Formula
size, memory, and runtime are therefore measured empirically in RQ2 rather than
inferred from asymptotic size alone.

\subsection{Implied constraints}

The primary treatment \texttt{both} combines two logically implied constraints
intended to strengthen propagation.  First, let $t_{sh}$ denote that service $s$
uses slot $h$, and let $v_{uh}$ denote that user $u$ uses $h$.  We channel
$t_{sh}\leftrightarrow\bigvee_a x_{ash}$ and
$v_{uh}\leftrightarrow\bigvee_{s\in\mathcal S_u}t_{sh}$.  Because the benchmark
service lists form a user partition, full coverage and the user-slot at-most-one
constraints imply
\begin{equation}
  \sum_{h\in\mathcal H}v_{uh}=|\mathcal S_u|,
  \qquad \forall u\in\mathcal U.
\end{equation}

Second, for each slot $h$ we construct a bipartite candidate graph between
caregivers and users.  Edge $(a,u)$ exists when caregiver $a$ can perform at
least one service in $\mathcal S_u$ at $h$.  The assignments made in one slot
form a matching because of the caregiver-slot and user-slot constraints;
therefore
\begin{equation}
  \sum_{s\in\mathcal S}t_{sh}\le \nu(G_h),
\end{equation}
where $\nu(G_h)$ is the maximum matching size computed once during formula
construction.  Both constraints are logical consequences of the hard model and
remove no feasible schedule.  The bundled \texttt{both-plus} treatment is
excluded because, beyond \texttt{both}, it adds equivalent service-level projection,
workload-cap tightening, and service-slot clustering; all preserve feasible
schedules.

\subsection{Detected equivalences and safe symmetry breaking}

We group slots only when their binary feasibility vectors over every
caregiver--service pair are identical.  We group two services only when each
has one well-defined user and continuity set and they share that user,
continuity set, caregiver reward vector, and caregiver--slot feasibility
vector.  Swapping two members of either class consequently preserves all hard constraints and all
three objective values.

For consecutive equivalent values, $\mathit{earlier}_i$ and
$\mathit{later}_j$ denote assigning the earlier and later value to ordered
positions $i$ and $j$; we impose
$\mathit{later}_j\rightarrow\bigvee_{i<j}\mathit{earlier}_i$.  Positions are
services for slot classes and ordered caregiver--slot pairs for service classes.
These clauses retain a canonical representative of every
detected orbit.  In contrast to
implied constraints, they may remove labeled schedules, but they preserve
satisfiability and an optimal solution because each removed schedule has an
objective-equivalent canonical representative~\cite{BogaertsEtAl2022}.
Regression tests cover both
cardinality encodings, full and soft coverage, satisfiable and infeasible cases,
and all symmetry modes; the campaign verifier separately reconstructs
feasibility and objective values from every reported assignment.

\section{Experimental Design}
\label{sec:experiments}

\subsection{Configurations, benchmarks, and campaign}

Baseline configuration $B$ uses sorting networks without implied constraints
or symmetry breaking.  Reference configuration $R$ combines Totalizer,
the \texttt{both} implied-constraint treatment, and symmetry breaking for
detected slot and service equivalences; the label does not assume a speedup.
Here $B$ denotes an encoding configuration and B0 the weighted policy in
Eq.~\eqref{eq:weighted}.  A weighted $2\times2\times2$ factorial on 80 original
instances (16 classes, seeds 1--5) crosses encoding, \texttt{both}, and symmetry
factors.  Only this factorial, rather than the bundled $B$--$R$ contrast,
supports factor-level attribution.

The end-to-end $B$--$R$ contrast reuses the corresponding two cells of this
factorial, yielding 80 paired instances without a duplicate campaign.  Policy
evaluation uses only $R$: B0 and LEX-COS are run on 140 original instances (14
classes, seeds 1--10) under the same 300-s budget.  The two classes inspected
during commercial development are excluded.  LEX-OCS
(OT${\to}$CONT${\to}$SIM) is evaluated as an order-sensitivity analysis on the
pre-specified 70-instance subset with seeds 1--5.

Validation uses 80 evaluation-critical corrected-v2 instances (16 strata,
five evaluation seeds per stratum) under B0 and LEX-COS with $R$.  A separate
20-instance commercial subset compares exact objectives from EvalMaxSAT, Gurobi
MIP, and CPLEX MIP under both policies; commercial runtimes are not used to
attribute MaxSAT encoding effects.  The $\varepsilon$-constraint, weight,
uncertainty, and routing branches are outside the measured compact campaign.

\subsection{Protocol, measures, and analysis}

Experiments run on one non-Spot GCP \texttt{c4-highcpu-8} VM (8 vCPUs, 16\,GB
RAM, Ubuntu 24.04 LTS).  We use the same EvalMaxSAT solver family selected by
the earlier study and pin the Linux x86-64 binary by SHA-256 in the artifact.
Every measured run uses one solver process restricted to one pinned vCPU, with
one campaign worker and no concurrent workload; Gurobi and CPLEX additionally
use one native solver thread.  The C++ encoder uses
\texttt{-O3 -DNDEBUG -std=c++11}.  Ten fixed warm-up instances from the
disjoint calibration set are unmeasured; a seeded, instance-major blocked order
interleaves configurations thereafter.  Every measured run has a 300\,s
end-to-end timeout; a lexicographic timeout is
cumulative across all stages.  Gurobi and CPLEX use one thread, seed~0, and zero
relative and absolute optimality gaps.  The earlier study used a 3600\,s limit
on a different processor, so we do not compare absolute runtimes across studies.

We record formula size, encoding/solver/end-to-end time, peak RSS, termination,
verifier outcome, and independently recomputed SIM, CONT, and OT.  Solved count
includes solver-reported \textsc{Optimum} and established infeasibility; PAR-2
assigns $2T$ to every other run under timeout $T$.  Objective deltas and
secondary runtime ratios use matched pairs only when both assignments are
optimality-reported and verifier-passing.  All eight factorial cells are
reported; the artifact retains all 12 direct factor contrasts, while the main
table displays the contrasts needed to expose encoding contexts and interaction
claims.  We stratify by instance class when aggregate effects reverse.  Timeouts
remain in solved-count and PAR-2 summaries.

The artifact retains source, instance, binary, and solver hashes; exact
commands; raw JSON and native logs; manifests; environment metadata; and
verifier status.  Pinned analysis scripts regenerate every result table and the
frozen-evidence marker.  Runtime claims are limited to this EvalMaxSAT and GCP
protocol.

\iffrozenresults
  % Generated only from the validated, frozen campaign artifact.
  % generated/results.tex  — EvalMaxSAT exploratory campaign
% DATA PROVENANCE — MANDATORY NOTE:
%   Solver: EvalMaxSAT (EXPLORATORY). Primary campaign requires a fresh frozen run.
%   Weighted factorial: 250 total instances, 176 feasible, class 30_10_4, 8 configs, T=300s.
%   LEX-OCS: 166 instances/config, 108 unique feasible, 8 configs.
%   LEX-COS: 45 instances/config, 0 OPTIMUM (all timeout with EvalMaxSAT).
%   Commercial: 94 instances, class 30_15_4, Gurobi + CPLEX, weighted + lex methods.
%   Replace this file with fresh publication-campaign outputs before submission.

\section{Results}
\label{sec:results}

\subsection{RQ2--RQ3: Encoding size and factorial effects}

Table~\ref{tab:factorial} reports solved counts and PAR-2 across all eight
cells of a $2{\times}2{\times}2$ factorial on 176 feasible instances
(250 total; 74 independently verified \textsc{Unsatisfiable}).
All runs used EvalMaxSAT with a 300-second timeout.

\begin{table}[t]
  \centering
  \caption{Eight-cell factorial results (weighted, EvalMaxSAT,
  $n=176$ feasible, $T=300$\,s).
  \textsc{Opt}: solver-reported optimum; PAR-2: $2T$ for timeouts.
  Card.\ = cardinality encoding; IC = implied constraints;
  SB = symmetry breaking.
  \textit{Exploratory; not publication evidence.}}
  \label{tab:factorial}
  \scriptsize
  \begin{tabular}{@{}lccrrr@{}}
    \toprule
    Configuration & Card. & IC & SB & \textsc{Opt} & PAR-2 (s) \\
    \midrule
    ORIGINAL      & SN  & \textemdash & \textemdash & 161 & 100.0 \\
    SN-none-ss    & SN  & \textemdash & ss           & 160 & 102.0 \\
    SN-both-none  & SN  & both        & \textemdash & 158 & 108.2 \\
    SN-both-ss    & SN  & both        & ss           & 159 & 108.7 \\
    \midrule
    TOT-none-none & TOT & \textemdash & \textemdash & 163 &  90.4 \\
    TOT-none-ss   & TOT & \textemdash & ss           & 163 &  91.7 \\
    TOT-both-none & TOT & both        & \textemdash & 162 &  96.6 \\
    TOT-both-ss   & TOT & both        & ss           & 162 &  96.3 \\
    \bottomrule
  \end{tabular}
\end{table}

\paragraph{Encoding (RQ2).}
Switching from sorting networks to Totalizer reduces the median variable count
from 57,518 to 20,650 ($-$64\%); the hard-clause count increases from 222,629
to 263,634 ($+$19\%), while the soft-clause count is identical (median 4,081).
On the 161 instances solved under both ORIGINAL and TOT-none-none, Totalizer
is faster on 113 and slower on 48 (median paired speedup~1.12; TOT-none-none:
163~solved, PAR-2 90.4\,s vs.\ ORIGINAL: 161~solved, 100.0\,s).

\paragraph{Implied constraints (RQ3).}
Adding \texttt{both} to a sorting-network base (SN-both-none vs.\ ORIGINAL)
raises PAR-2 from 100.0 to 108.2\,s and lowers the solved count by~3.
On 156 matched pairs, \texttt{both} is slower on 115 and faster on 41
(median ratio 0.90, i.e., $\approx$10\% slower).
On a Totalizer base the cost is smaller (TOT-both-none: PAR-2 96.6\,s vs.\
TOT-none-none: 90.4\,s, $-$1 solved), indicating a negative encoding
interaction whose magnitude depends on the cardinality base.

\paragraph{Symmetry breaking (RQ3).}
Slot-service symmetry breaking (ss) has a small, non-monotone effect:
PAR-2 increases by 1.3--1.7\,s on SN and 1.3\,s on TOT, with no consistent
direction of change in solved count.  Stratification by instance class is
needed before attributing the effect to symmetry breaking alone.

\subsection{RQ1: Weighted and lexicographic policies}

\paragraph{LEX-OCS order-sensitivity (exploratory).}
Across eight configurations on 108 unique feasible instances (class
\texttt{30\_10\_4}), the LEX-OCS policy (OT$\to$CONT$\to$SIM) reaches
\textsc{Optimum} on 841 of 849 feasible runs (99\%),
with at most~1 timeout per configuration.
On the 208 run-pairs solved to \textsc{Optimum} under both B0 and LEX-OCS
(same configuration and instance), LEX-OCS yields median
$\Delta\mathrm{CONT}=-4$ (all 208 pairs strictly negative),
$\Delta\mathrm{OT}=0$ (all 208 pairs zero), and
$\Delta\mathrm{SIM}=-12$ (all 208 pairs strictly negative).
Optimizing OT first eliminates overtime on this benchmark but reduces the
similarity score, reflecting the expected policy trade-off.

\paragraph{LEX-COS (EvalMaxSAT, exploratory).}
The LEX-COS policy produced 0~\textsc{Optimum} completions across
45~instances and 8~configurations (192 feasible runs) under EvalMaxSAT
within 300\,s; all feasible runs timed out.
This reflects one incomplete historical EvalMaxSAT campaign and does not
characterize LEX-COS performance under the fresh, pre-specified publication
protocol.

\subsection{Cross-solver validation}

On 94 instances (class \texttt{30\_15\_4}) solved to \textsc{Optimum} by both
Gurobi and CPLEX under the weighted MIP formulation, the two backends agree on
all 94~weighted scores.  On 19 instances overlapping with the EvalMaxSAT pilot,
EvalMaxSAT also agrees on all 19~scores, confirming cross-paradigm consistency
for the weighted policy.

On the same 94~Gurobi-optimal instances, Gurobi's LEX-COS (lex-continuity)
solution lowers CONT in 91 of 94~pairs (3~ties, 0~worse; median
$\Delta\mathrm{CONT}=-4$).  OT increases in 7~pairs and is unchanged in~87.
SIM decreases in 91 pairs (median $\Delta\mathrm{SIM}=-10$).
This pattern is consistent with the continuity-first policy on a benchmark with
sparse excess workload.

\outlineblock{Results above are historical EvalMaxSAT exploratory data and
Gurobi/CPLEX development data on class 30\_15\_4.  Before submission replace
them with the fresh 732-row EvalMaxSAT/Gurobi/CPLEX publication campaign and
update all numbers from the frozen artifact.}

\else
  \ifinternaloutline
    % Internal-only layout contract for the generated Results section.
% No value in this file is publication evidence.  The frozen generator must
% reproduce this structure from validated EvalMaxSAT/Gurobi/CPLEX artifacts.

\section{Results}
\label{sec:results}

\outlineblock{Open with the campaign size and state in one sentence that all
reported optimal assignments pass verification.}

\subsection{RQ1: Objective-rule effects}
\label{sec:rq1}

\outlineblock{Lead with the weighted objective versus the continuity-first
objective under the enhanced encoding on the 42 original
instances.  Both objective rules use the enhanced encoding and
the same 300-s budget.  Report solved counts and
PAR-2 on all runs; compute objective deltas only on matched, verifier-passing
pairs where both runs report OPTIMUM.  Then report overtime-first sensitivity
on the 48 corrected-benchmark instances, where overtime variation is stronger.}

\paragraph{Weighted versus continuity-first objective.}
Across \resultplaceholder{n} matched OPTIMUM pairs, the continuity-first
objective changes median additional-caregiver count (CONT) by
\resultplaceholder{x}, overtime workload (OT) by \resultplaceholder{x}, and
compatibility score (SIM) by \resultplaceholder{x}.  The corresponding
improved/equal/worse
counts and the all-run PAR-2 values are reported in
Table~\ref{tab:policy-validation}A.

\paragraph{Priority-order sensitivity.}
The continuity-first and overtime-first objectives return the same verified
objective vector on
\resultplaceholder{n}/\resultplaceholder{n} matched OPTIMUM pairs.  On the
remaining pairs, Table~\ref{tab:policy-validation}A reports the direction and
magnitude of each objective change.

\subsection{RQ2: Totalizer encoding}
\label{sec:rq2}

\outlineblock{Use Table~\ref{tab:factorial}A for all eight cells and
Table~\ref{tab:factorial}B for paired factor comparisons.  In the paragraph report
formula structure, peak RSS, proved count, PAR-2, and the paired effect with its
bootstrap confidence interval.  A speedup above one favors the second configuration.
If the interval crosses one or class effects reverse, describe the result as
mixed rather than faster.}

Relative to sorting networks, Totalizer changes median variables by
\resultplaceholder{x}, hard clauses by \resultplaceholder{x}, and soft clauses
by \resultplaceholder{x}.  Its median runtime ratio is
\resultplaceholder{x} (95\% bootstrap confidence interval
\resultplaceholder{lo--hi}) over \resultplaceholder{n} paired runs for which
both configurations prove optimality.

\subsection{RQ3: Implied constraints and symmetry breaking}
\label{sec:rq3}

\outlineblock{Read the direct implied-constraint and symmetry-breaking comparisons in
Table~\ref{tab:factorial}B, not the bundled baseline--enhanced contrast.  Report
their interactions with the counting encoding
interactions when their direction differs.  Negative or neutral effects are
first-class findings.  Do not attribute an aggregate difference to one factor
unless its paired comparison supports the claim.}

The implied-constraint treatment has paired ratio
\resultplaceholder{x} under sorting networks and \resultplaceholder{x} under
Totalizer.  Symmetry breaking for interchangeable slots and services has paired ratio
\resultplaceholder{x}; its win/tie/loss counts are
\resultplaceholder{w/t/l}.

\begin{table*}[t]
  \caption{Factorial and end-to-end encoding results on 48 original instances
  ($T=300$\,s).  Panel~A reports all eight configurations; Panel~B reports the
  direct factor and baseline--enhanced comparisons.  SN denotes sorting network,
  TOT Totalizer, IC implied constraints, SB symmetry breaking, W/T/L
  wins/ties/losses, and CI confidence interval.  \textsc{Opt}, \textsc{Unsat},
  and TO denote optimal, infeasible, and timeout runs.  Speedup is left runtime
  divided by right runtime; PAR-2 and RSS are lower-is-better.}
  \label{tab:factorial}
  \centering
  \scriptsize
  \setlength{\tabcolsep}{3.5pt}
  \renewcommand{\arraystretch}{0.94}
  \begin{tabularx}{\textwidth}{@{}Xlllrrrrr@{}}
    \toprule
    \multicolumn{9}{@{}l}{\textit{Panel A: factorial cells and all-run performance}} \\
    Configuration & Counting & IC & SB & \textsc{Opt} & \textsc{Unsat} & TO & PAR-2 (s) & RSS (MB) \\
    \midrule
    SN-off-off  & SN  & off & off & \resultplaceholder{n} & \resultplaceholder{n} & \resultplaceholder{n} & \resultplaceholder{x} & \resultplaceholder{x} \\
    SN-off-on   & SN  & off & on  & \resultplaceholder{n} & \resultplaceholder{n} & \resultplaceholder{n} & \resultplaceholder{x} & \resultplaceholder{x} \\
    SN-on-off   & SN  & on  & off & \resultplaceholder{n} & \resultplaceholder{n} & \resultplaceholder{n} & \resultplaceholder{x} & \resultplaceholder{x} \\
    SN-on-on    & SN  & on  & on  & \resultplaceholder{n} & \resultplaceholder{n} & \resultplaceholder{n} & \resultplaceholder{x} & \resultplaceholder{x} \\
    TOT-off-off & TOT & off & off & \resultplaceholder{n} & \resultplaceholder{n} & \resultplaceholder{n} & \resultplaceholder{x} & \resultplaceholder{x} \\
    TOT-off-on  & TOT & off & on  & \resultplaceholder{n} & \resultplaceholder{n} & \resultplaceholder{n} & \resultplaceholder{x} & \resultplaceholder{x} \\
    TOT-on-off  & TOT & on  & off & \resultplaceholder{n} & \resultplaceholder{n} & \resultplaceholder{n} & \resultplaceholder{x} & \resultplaceholder{x} \\
    TOT-on-on   & TOT & on  & on  & \resultplaceholder{n} & \resultplaceholder{n} & \resultplaceholder{n} & \resultplaceholder{x} & \resultplaceholder{x} \\
    \midrule
    \multicolumn{9}{@{}l}{\textit{Panel B: paired factor comparisons and baseline--enhanced validation}} \\
    Contrast (left $\to$ right) & \multicolumn{2}{c}{Condition} & Pairs & W/T/L & \multicolumn{2}{c}{Median speedup [95\% CI]} & $\Delta$variables & $\Delta$hard clauses \\
    \midrule
    SN $\to$ TOT & \multicolumn{2}{l}{IC off, SB off} & \resultplaceholder{n} & \resultplaceholder{w/t/l} & \multicolumn{2}{c}{\resultplaceholder{x [lo, hi]}} & \resultplaceholder{x} & \resultplaceholder{x} \\
    SN $\to$ TOT & \multicolumn{2}{l}{IC on, SB on} & \resultplaceholder{n} & \resultplaceholder{w/t/l} & \multicolumn{2}{c}{\resultplaceholder{x [lo, hi]}} & \resultplaceholder{x} & \resultplaceholder{x} \\
    IC off $\to$ on & \multicolumn{2}{l}{TOT, SB off} & \resultplaceholder{n} & \resultplaceholder{w/t/l} & \multicolumn{2}{c}{\resultplaceholder{x [lo, hi]}} & \resultplaceholder{x} & \resultplaceholder{x} \\
    SB off $\to$ on & \multicolumn{2}{l}{TOT, IC on} & \resultplaceholder{n} & \resultplaceholder{w/t/l} & \multicolumn{2}{c}{\resultplaceholder{x [lo, hi]}} & \resultplaceholder{x} & \resultplaceholder{x} \\
    Baseline $\to$ enhanced & \multicolumn{2}{l}{same factorial, $n=48$} & \resultplaceholder{n} & \resultplaceholder{w/t/l} & \multicolumn{2}{c}{\resultplaceholder{x [lo, hi]}} & -- & -- \\
    \bottomrule
  \end{tabularx}
\end{table*}

\subsection{Validation and scope}
\label{sec:validation}

\outlineblock{Panel B tests whether the directional objective findings persist on
the 48 corrected-benchmark instances, whose overtime
pressure is stronger.  Panel C reports objective agreement only when all three
solvers prove optimum.  Do not compare commercial runtime with MaxSAT runtime.
The compact campaign requires both validation branches; missing rows stop the
freeze instead of silently reducing the scope.}

On the corrected benchmark, \resultplaceholder{directional finding}.  Across
the pre-specified commercial subset, all three exact solvers prove optimum on
\resultplaceholder{n} weighted and \resultplaceholder{n} continuity-first
groups, with
\resultplaceholder{n} objective disagreements.

\begin{table*}[t]
  \caption{Objective-rule and cross-solver results.  Panels~A and B compare
  objective rules on the original and corrected benchmarks; deltas are second
  rule minus first, with lower CONT/OT and higher SIM preferred.  Panel~C
  reports agreement where EvalMaxSAT, Gurobi, and CPLEX all prove optimality.
  PAR-2 includes timeouts and is lower-is-better.}
  \label{tab:policy-validation}
  \centering
  \scriptsize
  \setlength{\tabcolsep}{3.5pt}
  \renewcommand{\arraystretch}{0.94}
  \begin{tabularx}{\textwidth}{@{}Xlrrrrrrr@{}}
    \toprule
    \multicolumn{9}{@{}l}{\textit{Panel A: paired objective-rule effects}} \\
    Comparison & Encoding & $n$ & Proved 1/2 & Both optimal & $\Delta$SIM & $\Delta$CONT & $\Delta$OT & PAR-2 1/2 (s) \\
    \midrule
    Weighted $\to$ continuity-first (original) & Enhanced & 42 & \resultplaceholder{n/n} & \resultplaceholder{n} & \resultplaceholder{x} & \resultplaceholder{x} & \resultplaceholder{x} & \resultplaceholder{x/x} \\
    Continuity-first $\to$ overtime-first (corrected) & Enhanced & 48 & \resultplaceholder{n/n} & \resultplaceholder{n} & \resultplaceholder{x} & \resultplaceholder{x} & \resultplaceholder{x} & \resultplaceholder{x/x} \\
    \midrule
    \multicolumn{9}{@{}l}{\textit{Panel B: corrected-benchmark confirmatory set, enhanced encoding}} \\
    Objective rule & Encoding & Runs & \textsc{Opt} & TO & Median SIM & Median CONT & Median OT & PAR-2 (s) \\
    \midrule
    Weighted & Enhanced & 48 & \resultplaceholder{n} & \resultplaceholder{n} & \resultplaceholder{x} & \resultplaceholder{x} & \resultplaceholder{x} & \resultplaceholder{x} \\
    Continuity-first & Enhanced & 48 & \resultplaceholder{n} & \resultplaceholder{n} & \resultplaceholder{x} & \resultplaceholder{x} & \resultplaceholder{x} & \resultplaceholder{x} \\
    \midrule
    \multicolumn{9}{@{}l}{\textit{Panel C: objective agreement on the 20-instance commercial subset}} \\
    Three-solver comparison & Objective rule & Groups & All optimal & Agree & Disagree & \multicolumn{3}{c}{Compared measures} \\
    \midrule
    EvalMaxSAT / Gurobi / CPLEX & Weighted & 20 & \resultplaceholder{n} & \resultplaceholder{n} & \resultplaceholder{n} & \multicolumn{3}{c}{coverage, weighted score} \\
    EvalMaxSAT / Gurobi / CPLEX & Continuity-first & 20 & \resultplaceholder{n} & \resultplaceholder{n} & \resultplaceholder{n} & \multicolumn{3}{c}{coverage, SIM, CONT, OT} \\
    \bottomrule
  \end{tabularx}
\end{table*}

\paragraph{Excluded extensions.}
\outlineblock{State in one sentence that other exploratory objective,
uncertainty, and routing branches are outside the measured compact campaign.
Do not use development or imported exploratory rows as publication evidence.}

  \fi
\fi

\section{Threats to Validity}
\label{sec:threats}

The original benchmark contains limited excess-workload variation, so it may
understate differences involving OT.\@ Corrected-v2 supplies feasible witnesses
and stronger excess-workload pressure, but remains synthetic rather than
operational data.  Its disjoint calibration and evaluation sets limit tuning
leakage, not this external-validity threat.  Routing and uncertain availability
are outside the implemented decision model, and no claim about either extension
follows from the present experiments.

LEX-COS restarts EvalMaxSAT at each stage and shares one cumulative time budget.
Thus its runtime combines objective policy and staged implementation, and a
partial prefix is never counted as a lexicographic optimum.  Runtime evidence
also comes from one pinned MaxSAT solver version.  The Linux EvalMaxSAT binary
must pass weighted and three-stage LEX-COS WCNF smoke tests before any measured
run; this controls compatibility but does not establish solver generality.
Gurobi and CPLEX solve
the same exact MIP formulation to check objective agreement across paradigms on
a pre-specified 20-instance subset; that formulation differs from the MaxSAT
encoding, so they are validation backends, not a basis for claiming
solver-independent MaxSAT speedups.

Timeout censoring and class heterogeneity can distort solved-only averages.
We therefore retain timeouts in solved counts and PAR-2, restrict paired time
ratios to mutually optimal, verifier-passing cases, and stratify reversals by
instance class.  Finally, $B$ is intentionally an unstrengthened baseline
comparator, while $R$
bundles three changes.  The paired factorial, rather than that bundled contrast,
supports statements about individual factors.  Independent solution
verification, held-out instances, raw logs, and cryptographic hashes mitigate
implementation and provenance errors but cannot eliminate the scope limitations
above.  The verifier checks feasibility, inherited bounds, and objective values;
it does not check an independently generated MaxSAT proof trace, so optimality
labels remain backend-reported.

% Keep the two full-width Results tables ahead of the Conclusion heading.  In
% double-column mode they otherwise may leapfrog a short conclusion while
% waiting for the next top-of-page float slot.
\FloatBarrier
\iffrozenresults
  % generated/conclusion.tex — EvalMaxSAT exploratory campaign
% DATA PROVENANCE — MANDATORY NOTE:
%   Claims here are limited to EvalMaxSAT exploratory results.
%   Replace with fresh frozen publication-campaign conclusions before submission.

\section{Conclusion}
\label{sec:conclusion}

We studied lexicographic MaxSAT for home-care resource allocation by combining
the LEX-COS policy with bidirectional Totalizer encodings, logically implied
constraints, and symmetry-breaking for detected slot and service equivalences.

\paragraph{RQ2 (Totalizer encoding).}
On 176 feasible instances (EvalMaxSAT, class
\texttt{30\char`\_10\char`\_4}), the Totalizer encoding reduces the variable
count by 64\% relative to sorting networks and solves 2~additional instances
under a 300-second budget (163 vs.\ 161), with PAR-2 improving from 100.0 to
90.4\,s and a median paired speedup of~1.12.
The formula size reduction does not uniformly accelerate solving---the
hard-clause count increases 19\%---confirming that asymptotic size alone
does not predict runtime and motivating empirical evaluation.

\paragraph{RQ3 (Implied constraints and symmetry breaking).}
Adding implied constraints to a sorting-network base increases PAR-2 by 8\,s
and reduces solved count by 3, indicating a negative encoding interaction on
this benchmark.  On a Totalizer base the penalty is smaller.  Symmetry breaking
has a small mixed effect ($<$2\,s) without a clear monotone pattern.

\paragraph{RQ1 (Objective policies, exploratory).}
Gurobi's LEX-COS results on 94 instances confirm the expected policy behavior:
CONT decreases in 91 of 94 pairs relative to weighted (median $\Delta=-4$),
OT is unchanged in 87 pairs, and SIM decreases in 91 pairs (median $\Delta=-10$).
Cross-paradigm agreement on weighted scores is perfect across 94 Gurobi--CPLEX
pairs and 19 EvalMaxSAT--Gurobi pairs.  EvalMaxSAT timed out on all LEX-COS
instances; this historical incomplete run is not publication evidence.

\paragraph{Limitations and future work.}
The reported factorial results use EvalMaxSAT on class \texttt{30\_10\_4} and
are preliminary; the fresh EvalMaxSAT publication campaign on the pre-specified
compact matrix remains to be completed.  Routing and uncertain availability are outside the
implemented model and are natural next steps.  Native lexicographic solving and
validation on operational data would strengthen both generality and practitioner
relevance.

\outlineblock{Replace this conclusion with one derived from frozen publication
campaign outputs.  Claims about LEX-COS performance require the fresh EvalMaxSAT
data under the locked hash, sample, timeout, and provenance contract.}

\else
  \ifinternaloutline
    \section{Conclusion}
    \label{sec:conclusion}
    \outlineblock{Target 0.15--0.20 page.  Answer RQ1--RQ3 in result order using
    only frozen evidence.  Report mixed or conditional encoding effects as such.
    Close with incremental/native lexicographic solving and validation on
    operational data; routing and uncertainty remain model extensions.}
  \fi
\fi

% Artifact statement to activate after release creation:
% \section*{Data and Code Availability}
% Source, instances, manifests, raw results, and checksums are archived at
% \url{PERSISTENT-ARTIFACT-URL} for commit/tag \texttt{FROZEN-TAG}.

\bibliographystyle{ACM-Reference-Format}
\bibliography{references}

\end{document}
